\documentclass[a4paper,10pt]{report}\input{../0.Préambule/Préambule_cours_prof}\begin{document}

\pagestyle{empty}
\newpage
\subsection*{Développement de la pyramide}
\addcontentsline{toc}{subsection}{Développement de la pyramide}

On a vu ensemble dans le cours tous les développement du cube.  On en a trouvé $11$, tous différents (que l'on ne peux pas superposer exactement en les déplaçant ou en les retournant) et on a vérifié qu'il n'y en a pas d'autres.

\medskip
\noindent \begin{minipage}{10cm}
Aujourd'hui, je souhaiterai savoir combien il existe de développement d'une pyramide régulière de base carrée dont les quatre faces latérales son des triangles équilatéraux et dont toutes les arêtes mesurent $2$cm.

\medskip
J'en ai déjà trouvé $2$, mais je pense qu'il y en a encore d'autres :
\end{minipage}
\hspace{5mm}
\begin{minipage}{6cm}
\begin{center}
\includegraphics[scale=0.4]{../40.Image/16. Act-Espace - Pavé droit et volumes 1.png}
\end{center}
\end{minipage}

\medskip
\noindent \textbf{Combien en existe-t-il en tout, de développement différents de cette pyramide ?}

\medskip
\noindent \textbf{Dessinez les tous, avec tous les cotés de 2cm.}

\bigbreak
\subsection*{Développement de la pyramide}

On a vu ensemble dans le cours tous les développement du cube. On en a trouvé $11$, tous différents (que l'on ne peux pas superposer exactement en les déplaçant ou en les retournant) et on a vérifié qu'il n'y en a pas d'autres.

\medskip
\noindent \begin{minipage}{10cm}
Aujourd'hui, je souhaiterai savoir combien il existe de développement d'une pyramide régulière de base carrée dont les quatre faces latérales son des triangles équilatéraux et dont toutes les arêtes mesurent $2$cm.

\medskip
J'en ai déjà trouvé $2$, mais je pense qu'il y en a encore d'autres :
\end{minipage}
\hspace{5mm}
\begin{minipage}{6cm}
\begin{center}
\includegraphics[scale=0.4]{../40.Image/16. Act-Espace - Pavé droit et volumes 1.png}
\end{center}
\end{minipage}

\medskip
\noindent \textbf{Combien en existe-t-il en tout, de développement différents de cette pyramide ?}

\medskip
\noindent \textbf{Dessinez les tous, avec tous les cotés de 2cm.}

\bigbreak
\subsection*{Développement de la pyramide}

On a vu ensemble dans le cours tous les développement du cube. On en a trouvé $11$, tous différents (que l'on ne peux pas superposer exactement en les déplaçant ou en les retournant) et on a vérifié qu'il n'y en a pas d'autres.

\medskip
\noindent \begin{minipage}{10cm}
Aujourd'hui, je souhaiterai savoir combien il existe de développement d'une pyramide régulière de base carrée dont les quatre faces latérales son des triangles équilatéraux et dont toutes les arêtes mesurent $2$cm.

\medskip
J'en ai déjà trouvé $2$, mais je pense qu'il y en a encore d'autres :
\end{minipage}
\hspace{5mm}
\begin{minipage}{6cm}
\begin{center}
\includegraphics[scale=0.4]{../40.Image/16. Act-Espace - Pavé droit et volumes 1.png}
\end{center}
\end{minipage}

\medskip
\noindent \textbf{Combien en existe-t-il en tout, de développement différents de cette pyramide ?}

\medskip
\noindent \textbf{Dessinez les tous, avec tous les cotés de 2cm.}

\bigbreak
\subsection*{Développement de la pyramide}

On a vu ensemble dans le cours tous les développement du cube. On en a trouvé $11$, tous différents (que l'on ne peux pas superposer exactement en les déplaçant ou en les retournant) et on a vérifié qu'il n'y en a pas d'autres.

\medskip
\noindent \begin{minipage}{10cm}
Aujourd'hui, je souhaiterai savoir combien il existe de développement d'une pyramide régulière de base carrée dont les quatre faces latérales son des triangles équilatéraux et dont toutes les arêtes mesurent $2$cm.

\medskip
J'en ai déjà trouvé $2$, mais je pense qu'il y en a encore d'autres :
\end{minipage}
\hspace{5mm}
\begin{minipage}{6cm}
\begin{center}
\includegraphics[scale=0.4]{../40.Image/16. Act-Espace - Pavé droit et volumes 1.png}
\end{center}
\end{minipage}

\medskip
\noindent \textbf{Combien en existe-t-il en tout, de développement différents de cette pyramide ?}

\medskip
\noindent \textbf{Dessinez les tous, avec tous les cotés de 2cm.}

\newpage
\subsection*{Développement de la pyramide - Feuille du professeur}
\addcontentsline{toc}{subsection}{Développement de la pyramide - Feuille du professeur}

\paragraph{Tâche de résolution et savoirs mobilisés}

\begin{enumerate}
\item[•]Vérifier que les deux exemples de développements permettent bien de construire la pyramide. 
\item[•]Organiser la recherche pour obtenir les autres développements en esquissant des schémas. 

Par exemple, placer la base et disposer les triangles de manière systématique : 

\begin{enumerate}
\item[•]un par côté (exemple donné) « 1, 1, 1, 1 » ; 
\item[•]répartis sur trois côtés « 0, 1, 1, 2 », 
\item[•]répartis sur 2 côtés adjacents ou opposés, 
\item[•]par groupes de deux : « 0, 0, 2, 2 » ; « 0, 2, 0, 2 » ; 
\item[•]répartis sur 2 côtés adjacents ou opposés avec un groupe de trois et un isolé : « 0, 0, 1, 3 » ; « 0, 1, 0, 3 », 
\item[•]tous attachés à un côté : « 0, 0, 0, 4 ». 
\end{enumerate}
\item[•]Essayer les différentes possibilités pour chacune des dispositions précédentes et éliminer celles qui donnent des faces superposées après le pliage. 
\item[•]Éliminer les figures isométriques (superposables par un déplacement dans le plan - translation rotation – ou par une symétrie axiale ou retournement) pour ne retenir que les huit développements différents de la figure suivante : 

\begin{center}
\includegraphics[scale=1]{../40.Image/16. Act-Espace - Pavé droit et volumes 2.png}
\end{center} 
 
\item[•]Après vérifications des schémas, effectuer les constructions aux dimensions requises (2 cm par côté) 

\item[]Ou : découper des carrés et des triangles équilatéraux de 2 cm d côté, construire la pyramide en collant les pièces puis « l'ouvrir » successivement de différentes manières et noter les dispositions des faces dans le plan avant d'éliminer celles qui sont isométriques. 
\end{enumerate}

\paragraph{Notation :}

Selon les critères déterminés lors de l'analyse a priori : 
\begin{enumerate}
\item[•]4 points: Réponse correcte (8 développements différents  avec les 8 dessins), on n'exige pas une grande précision, mais la distinction claire du carré, des quatre triangles équilatéraux et de leurs positions respectives
\item[•]3 points: Réponse correcte (8 développements différents) avec dessin des six développements qui ne sont pas donnés dans les exemples (sans doublures) 
ou réponse avec 7 ou 9 développements différents (due à un oubli ou une « doublure ») avec les dessins correspondants clairs (avec ou sans les exemples) 
\item[•]2 points: Réponse avec 5 ou 6 ou 10 ou 11 développements différents avec les dessins correspondants clairs, due à deux oublis ou deux « doublures » 
ou 8 développements différents mais avec dessins très imprécis ou faces non-contiguës
\item[•]1 point: De un à quatre développements trouvés, différents des exemples
\item[•]0 point: Incompréhension du problème
\end{enumerate}
\end{document}